\section{Training-arrival review and synchronized figures}
The report and separate comparison use the same 20 generated figure files and
63 panel records: 53 partial and ten unreproduced. The document builder now
records each panel's status, numerical-file hashes, shared figure hash, and PDF
file timestamps in \path{report/document_build.json}. All 19 comparison pages
were visually inspected. Published images remain confined to that comparison;
missing results retain explicit explanations. This synchronization does not
validate the plotted scientific results or promote the failed velocity profiles.

\paragraph{Targeted observations and review.}
Existing target-station files were checked before querying EIDA for
13 December 2024, 00:00--02:00 UTC. All six station requests succeeded, adding
36 core channel-hours and 20,906,496 bytes against a 20,519,430-byte estimate.
The 40 GiB free-space reserve plus 8 GiB processing allowance was preserved.
The broad queue remained stopped. Header epochs, responses, coordinates and
sampling rates agree with existing StationXML; absolute clock accuracy is untested.
The permitted trained PhaseNet weights yield 163 P and 922 S picks, but no
GaMMA association. CUMA has no picks in this window. No new arrival therefore
enters calibration. CGAU P/S, CLAC S, CREM P/S and CRUS P/S still have one
training observation each, from the preceding December 12 event. CSOC P/S
and CUMA P/S remain absent. Available public intervals beyond this bounded
window have not been exhausted or declared inaccessible.

Six inherited training arrivals remain ambiguous in enlarged raw and filtered
windows: CFMN P for event 159, CMSN S for 265, CFMN S for 264, and CSFT P,
CPOZ S and CFMN S for 1055. They are preserved as uncertain and excluded from
the new fit; no times were forced. Network review of event 1532 places several
assigned P markers in an earlier decaying packet, without establishing a distinct
new event. That association is also excluded as uncertain, not declared false.
After the unchanged event screen, 51 events and 816 arrivals remain; event 159
fails after its ambiguous P is removed. The selection sensitivity retains
62 events and 997 arrivals. These are Codex visual reviews, not independent
human seismologist validation. Other inherited arrivals were not all reviewed.

\paragraph{Numerical diagnosis and refit.}
Saved 32/64-bit prelocation runs have identical observation, model, station and
control files. Their first iteration-zero RMS values differ by about
$2\times10^{-6}$ s. The three-layer prelocation then differs by 0.11 km for
event 932 near the 3 km interface; the four-layer path reaches a 0.18 km
difference after its first joint block and 0.65 km later. This is arithmetic
sensitivity amplified along nonlinear paths, rather than a detected input mismatch.
The exact ray-branch cause has not been isolated for these updated cohorts.
An additional native logging defect leaves the aggregate P/S counters
uninitialized. Those variables are used only for printing; input tokens and
per-ray records supply the verified counts. No solver equation or binary was
changed to hide this defect.

All 18 training fits were frozen before evaluating the unchanged 62-event,
1,065-arrival 2025 set. The same three parameterizations and six sensitivity
choices were retained. No 2025 data were repicked, reassociated, reweighted or
used to fit velocities/delays or select a model. The existing evaluation adjusts
held-out origins and locations with velocities and delays fixed; it tests
conditional arrival prediction, not independent known-hypocenter accuracy.
\begin{center}\small\begin{tabular}{lccc}\toprule
Parameterization & Training RMS (s) & 2025 RMS (s) & All gates\\\midrule
Four layers, fitted delays & 0.2039 $\to$ 0.1584 & 0.2406 $\to$ 0.2432 & Fail \\
Three layers, fitted delays & 0.2044 $\to$ 0.1603 & 0.2432 $\to$ 0.2446 & Fail \\
Three layers, zero delays & 0.2044 $\to$ 0.1981 & 0.2432 $\to$ 0.2271 & Fail \\
\bottomrule\end{tabular}\end{center}
No candidate passes every unchanged criterion. The maximum depth changes across
the registered sensitivities are 1.65, 1.00, 1.09 km in table order.
The zero-delay case changes by 733 m horizontally and 0.32 km/s in velocity.
Additional training iterations in that case still move locations by 889 m horizontally and 0.56 km in depth.
Fixed-model held-out locations also fail the unchanged convergence check.
The limits remain 100 m horizontally, 0.1 km in depth, 0.05 km/s in velocity
and 0.02 s in station delay, alongside physical velocities, convergence and
improved holdout RMS. No failed result is promoted into the catalog or figures.

\paragraph{Depth support and what remains needed.}
In the zero-delay base fit, 53/363 P, 77/453 S rays extend strictly below 3 km.
These are model-dependent paths with 0.01 km vertex rounding, not measured depths
or independent proof of deep resolution. The December 12 event spans 2.70--3.06 km across the fits.
Independent layerwise travel-time sums agree with serialized residuals to
0.0050--0.0083 s RMS, subject to rounded native coordinates, velocities and times.
This is an arithmetic consistency check, not confirmation of phase identity.
The present observations do not establish a stable profile. Clear pre-2025
CSOC/CUMA arrivals recorded jointly with the established network are still
needed, as are repeat events at the seven singly represented station/phase
combinations. Training associations in coda and events sensitive to layer
interfaces require further scientific review. More independently constrained
deep source--receiver paths are needed to test resolution below 3 km; no exact
number of additional events sufficient for acceptance has been established.
Logs, decisions, fits, residuals and ray tables are in
\path{data/velocity_holdout_review/sync_training_review_20261015/}; commands are in
\path{docs/SYNC_VELOCITY_REVIEW.md}. Earlier runs remain separate checkpoints.
